% Slideaway — Thesis/Dissertation Defense Template (30–45 min)
% Theme: Boadilla/whale — formal, clean footline with author|title|page/total
% Structural scaffolding only. Style/color theory → design-foundations.md
% Engine workflow → engine-beamer.md

\documentclass[aspectratio=169,11pt]{beamer}

\usetheme{Boadilla}
\usecolortheme{whale}

\setbeamertemplate{navigation symbols}{}
% Custom footline: frame N/total for committee tracking
\setbeamertemplate{footline}{
  \leavevmode
  \hbox{
    \begin{beamercolorbox}[wd=.333333\paperwidth,ht=2.25ex,dp=1ex,center]{author in head/foot}
      \usebeamerfont{author in head/foot}\insertshortauthor
    \end{beamercolorbox}
    \begin{beamercolorbox}[wd=.333333\paperwidth,ht=2.25ex,dp=1ex,center]{title in head/foot}
      \usebeamerfont{title in head/foot}\insertshorttitle
    \end{beamercolorbox}
    \begin{beamercolorbox}[wd=.333333\paperwidth,ht=2.25ex,dp=1ex,right]{date in head/foot}
      \usebeamerfont{date in head/foot}
      \insertframenumber{}/\inserttotalframenumber\hspace*{2ex}
    \end{beamercolorbox}
  }
  \vskip0pt
}

% Section dividers
\AtBeginSection[]{
  \begin{frame}
    \vfill
    \centering
    \begin{beamercolorbox}[sep=8pt,center,shadow=true,rounded=true]{title}
      \usebeamerfont{title}\insertsectionhead\par
    \end{beamercolorbox}
    \vfill
  \end{frame}
}

% ── Packages ──────────────────────────────────────────────────────────────
\usepackage{booktabs}
\usepackage{graphicx}
\usepackage{amsmath}
\usepackage{hyperref}
\usepackage{xcolor}
\usepackage{algorithm}
\usepackage{algorithmic}
\usepackage{listings}
\usepackage{multirow}
\usepackage[backend=biber,style=authoryear]{biblatex}
% \addbibresource{references.bib}

\lstset{
  basicstyle=\ttfamily\scriptsize,
  keywordstyle=\color{blue},
  commentstyle=\color{gray},
  breaklines=true,
  frame=single,
}

% ── Title Metadata ────────────────────────────────────────────────────────
\title[Defense]{Dissertation Title: A Study of Something Important}
\subtitle{PhD Dissertation Defense}
\author[Candidate]{Candidate Name}
\institute{Department of X \\ University of Y}
\date{\today}

% ══════════════════════════════════════════════════════════════════════════
\begin{document}

\begin{frame}
  \titlepage
\end{frame}

% ── Committee ─────────────────────────────────────────────────────────────

\begin{frame}{Defense Committee}
  \begin{description}
    \item[Chair] Prof.\ A (Department of X)
    \item[Member] Prof.\ B (Department of Y)
    \item[Member] Prof.\ C (Department of Z)
    \item[External] Prof.\ D (Other University)
  \end{description}
  \note{Acknowledge each committee member by name. Set a respectful tone.}
\end{frame}

\begin{frame}{Outline}
  \tableofcontents
\end{frame}

% ── Dissertation Overview (2–3 slides, ~5 min) ───────────────────────────
\section{Dissertation Overview}

\begin{frame}{Motivation and Context}
  \begin{itemize}
    \item Broad problem area and its significance
    \item Specific gap in current knowledge
    \item Overarching research question
  \end{itemize}
  \note{Frame the entire dissertation in 90 seconds. Committee already read it — remind them of the arc.}
\end{frame}

\begin{frame}{Dissertation Structure}
  \begin{table}
    \caption{Chapter overview}
    \begin{tabular}{clc}
      \toprule
      Chapter & Topic & Status \\
      \midrule
      1 & Introduction & Background \\
      2 & Literature Review & Foundation \\
      3 & Study 1: [Title] & Published \\
      4 & Study 2: [Title] & Under Review \\
      5 & Study 3: [Title] & Complete \\
      6 & General Discussion & Synthesis \\
      \bottomrule
    \end{tabular}
  \end{table}
  \note{Show where each study fits. Publication status signals rigor.}
\end{frame}

% ── Study 1 (5–7 slides, ~8 min) ─────────────────────────────────────────
\section{Study 1: [Title]}

\begin{frame}{Study 1 — Motivation and Hypotheses}
  \begin{itemize}
    \item Specific question addressed by Study 1
    \item Hypothesis H1a: ...
    \item Hypothesis H1b: ...
  \end{itemize}
\end{frame}

\begin{frame}{Study 1 — Methods}
  \begin{columns}[T]
    \begin{column}{0.48\textwidth}
      \textbf{Data \& Design}
      \begin{itemize}
        \item Data source and coverage
        \item Experimental design
        \item Sample size and power
      \end{itemize}
    \end{column}
    \begin{column}{0.48\textwidth}
      \framebox[\textwidth]{\parbox{0.9\textwidth}{\centering\vspace{2cm}Study 1 Design\vspace{2cm}}}
    \end{column}
  \end{columns}
\end{frame}

\begin{frame}{Study 1 — Key Results}
  \begin{center}
    \framebox[0.8\textwidth]{\parbox{0.7\textwidth}{\centering\vspace{3cm}Study 1 Main Figure\vspace{3cm}}}
  \end{center}
  \note{Lead with the figure. State the key finding in one sentence.}
\end{frame}

\begin{frame}{Study 1 — Summary}
  \begin{itemize}
    \item H1a: Supported / Not supported — evidence summary
    \item H1b: Supported / Not supported — evidence summary
    \item Contribution to overall dissertation arc
  \end{itemize}
\end{frame}

% ── Study 2 (5–7 slides, ~8 min) ─────────────────────────────────────────
\section{Study 2: [Title]}

\begin{frame}{Study 2 — Motivation and Hypotheses}
  \begin{itemize}
    \item How Study 2 builds on Study 1 findings
    \item Hypothesis H2a: ...
    \item Hypothesis H2b: ...
  \end{itemize}
\end{frame}

\begin{frame}{Study 2 — Methods}
  \begin{columns}[T]
    \begin{column}{0.48\textwidth}
      \textbf{Data \& Design}
      \begin{itemize}
        \item Data source and coverage
        \item Design choices and rationale
        \item Analytical approach
      \end{itemize}
    \end{column}
    \begin{column}{0.48\textwidth}
      \framebox[\textwidth]{\parbox{0.9\textwidth}{\centering\vspace{2cm}Study 2 Design\vspace{2cm}}}
    \end{column}
  \end{columns}
\end{frame}

\begin{frame}{Study 2 — Key Results}
  \begin{center}
    \framebox[0.8\textwidth]{\parbox{0.7\textwidth}{\centering\vspace{3cm}Study 2 Main Figure\vspace{3cm}}}
  \end{center}
\end{frame}

\begin{frame}{Study 2 — Summary}
  \begin{itemize}
    \item H2a: Supported / Not supported
    \item H2b: Supported / Not supported
    \item Contribution to overall dissertation arc
  \end{itemize}
\end{frame}

% ── Study 3 (5–7 slides, ~8 min) ─────────────────────────────────────────
\section{Study 3: [Title]}

\begin{frame}{Study 3 — Motivation and Hypotheses}
  \begin{itemize}
    \item How Study 3 extends Studies 1 and 2
    \item Hypothesis H3a: ...
    \item Hypothesis H3b: ...
  \end{itemize}
\end{frame}

\begin{frame}{Study 3 — Methods}
  \begin{columns}[T]
    \begin{column}{0.48\textwidth}
      \textbf{Data \& Design}
      \begin{itemize}
        \item Data source and coverage
        \item Novel methodological contribution
        \item Validation approach
      \end{itemize}
    \end{column}
    \begin{column}{0.48\textwidth}
      \framebox[\textwidth]{\parbox{0.9\textwidth}{\centering\vspace{2cm}Study 3 Design\vspace{2cm}}}
    \end{column}
  \end{columns}
\end{frame}

\begin{frame}{Study 3 — Key Results}
  \begin{center}
    \framebox[0.8\textwidth]{\parbox{0.7\textwidth}{\centering\vspace{3cm}Study 3 Main Figure\vspace{3cm}}}
  \end{center}
\end{frame}

\begin{frame}{Study 3 — Summary}
  \begin{itemize}
    \item H3a: Supported / Not supported
    \item H3b: Supported / Not supported
    \item Contribution to overall dissertation arc
  \end{itemize}
\end{frame}

% ── General Discussion (3–4 slides, ~7 min) ───────────────────────────────
\section{General Discussion}

\begin{frame}{Cross-Study Synthesis}
  \begin{table}
    \caption{Hypothesis support across studies}
    \begin{tabular}{lccc}
      \toprule
      Hypothesis & Study 1 & Study 2 & Study 3 \\
      \midrule
      H1 & \checkmark & — & \checkmark \\
      H2 & — & \checkmark & \checkmark \\
      H3 & — & — & \checkmark \\
      \bottomrule
    \end{tabular}
  \end{table}
  \note{The synthesis table is the intellectual climax. Show how the pieces fit together.}
\end{frame}

\begin{frame}{Theoretical Contributions}
  \begin{enumerate}
    \item Contribution to theory or framework
    \item New understanding enabled by this work
    \item How this changes the field's perspective
  \end{enumerate}
\end{frame}

\begin{frame}{Practical Implications}
  \begin{itemize}
    \item Application 1: who benefits and how
    \item Application 2: policy or practice implications
    \item Broader impact statement
  \end{itemize}
\end{frame}

\begin{frame}{Limitations}
  \begin{itemize}
    \item Scope limitations across all three studies
    \item Methodological constraints and their impact
    \item Generalizability caveats
  \end{itemize}
  \note{Committee expects honest self-assessment. Address likely questions proactively.}
\end{frame}

% ── Future Research Program (1–2 slides) ──────────────────────────────────
\section{Future Research}

\begin{frame}{Future Research Program}
  \begin{enumerate}
    \item Immediate follow-up: address Study 3 limitations
    \item Next project: extend framework to new domain
    \item Long-term vision: research trajectory post-defense
  \end{enumerate}
  \note{Show you have a program, not just a dissertation. Important for committee evaluation.}
\end{frame}

% ── Conclusion ────────────────────────────────────────────────────────────
\section{Conclusion}

\begin{frame}{Dissertation Summary}
  \begin{enumerate}
    \item Study 1 established [finding]
    \item Study 2 demonstrated [finding]
    \item Study 3 revealed [finding]
    \item Together: [overarching contribution]
  \end{enumerate}
  \vfill
  \small
  Preprint: \url{https://arxiv.org/abs/xxxx.xxxxx} \\
  Code: \url{https://github.com/your/repo}
\end{frame}

\begin{frame}{Acknowledgments}
  \small
  This work was supported by [Funding Agency, Grant \#].

  Thank you to my advisor, committee, lab members, and family.

  % \printbibliography
\end{frame}

\begin{frame}[standout]
  Thank You — Questions?
\end{frame}

% ══════════════════════════════════════════════════════════════════════════
\appendix

\begin{frame}{Backup: Full Regression Tables}
  % Complete statistical output for all three studies
\end{frame}

\begin{frame}{Backup: Model Fit \& Diagnostics}
  % Residuals, convergence, goodness-of-fit across studies
\end{frame}

\begin{frame}{Backup: Alternative Analyses}
  % Robustness checks, sensitivity to assumptions
\end{frame}

\begin{frame}{Backup: Extended Literature}
  % Additional citations committee might ask about
\end{frame}

\begin{frame}{Backup: Methodological Details}
  % Implementation specifics, parameter choices, code architecture
\end{frame}

\end{document}
